Теоретическая постановка задачи JFO (уравнение сохранения~\eqref{eq:jfo_conservation}, условия комплементарности~\eqref{eq:jfo_complementarity}) и общая схема итерационного алгоритма active-set приведены в подразделах~2.3.2-2.3.3~\cite{elrod1975}. В настоящем разделе описывается численная реализация каждого шага этого алгоритма: решение подсистемы для давления на GPU (раздел~3.4), обновление степени заполнения~$\vartheta$ в кавитационной зоне, переключение узлов между зонами полного заполнения и кавитации, критерии остановки. Алгоритм JFO имеет двухуровневую итерационную структуру: внутренний цикл~-- решение линейной системы для давления методом Red-Black SOR на GPU при фиксированном~$\vartheta$, внешний цикл~-- обновление активного множества и пересчёт~$\vartheta$.

\textit{Решение подсистемы для давления.}
На каждой внешней итерации~$m$ фиксируется текущее распределение~$\vartheta^{(m)}$. В узлах зоны полного заполнения ($\vartheta = 1$) решается система~\eqref{eq:5point_stencil} методом Red-Black SOR (раздел~3.4) до сходимости по критерию~\eqref{eq:convergence_change}. В узлах зоны кавитации давление фиксировано: $p_{i,j} = p_{\mathrm{cav}}$. Эти узлы не обновляются решателем давления; фиксация обеспечивается проекцией $p \leftarrow \max(p,\, p_{\mathrm{cav}})$, встроенной в ядро обновления (раздел~3.4). В кавитационных узлах давление фиксируется на уровне~$p_{\mathrm{cav}}$, поэтому внутри связной кавитационной области градиенты давления отсутствуют и напорный поток не формируется. На границе кавитации и полного заполнения градиент давления, как правило, отличен от нуля и определяется решением в зоне полного заполнения.

Коэффициенты напорной части дискретизации ($a_E$, $a_W$, $a_N$, $a_S$, $a_P$), определяемые полем зазора~$h$, не зависят от~$\vartheta$ и могут переиспользоваться во внутреннем цикле. Зависимость от~$\vartheta$ проявляется в переносном (куэттовском) потоке~\eqref{eq:flux_couette}, а также в классификации узлов как кавитационных ($p = p_{\mathrm{cav}}$) или полностью заполненных.

\textit{Обновление степени заполнения в кавитационной зоне.}
В кавитационной зоне $p = p_{\mathrm{cav}}$; внутри связной кавитационной области градиенты давления отсутствуют и напорный поток не формируется. Уравнение JFO~\eqref{eq:jfo_conservation} сводится к переносу массы куэттовским потоком. Для стационарной задачи это условие баланса переносных потоков:
\begin{equation}\label{eq:jfo_couette_balance}
  \frac{F_{C,\,i+\frac{1}{2},j} - F_{C,\,i-\frac{1}{2},j}}{\Delta x_1}
  + \frac{G_{C,\,i,j+\frac{1}{2}} - G_{C,\,i,j-\frac{1}{2}}}{\Delta x_2}
  = 0,
\end{equation}
\noindent где переносные потоки на гранях определены в~\eqref{eq:flux_couette} (раздел~3.2):
\begin{equation}\label{eq:jfo_theta_flux}
  F_{C,\,i+\frac{1}{2},j} = \frac{\vartheta_{\mathrm{up}}\,h_{i+\frac{1}{2},j}}{2}\,U_1,
  \quad
  G_{C,\,i,j+\frac{1}{2}} = \frac{\vartheta_{\mathrm{up}}\,h_{i,j+\frac{1}{2}}}{2}\,U_2.
\end{equation}

Значение~$\vartheta_{\mathrm{up}}$ берётся из наветренного узла по направлению скорости~$\mathbf{U}$ (upwind-схема, раздел~3.2). Толщина на грани~-- арифметическое среднее значений в соседних узлах.

В двумерной постановке соотношение~\eqref{eq:jfo_couette_balance} рассматривается как дискретное уравнение переноса для~$\vartheta$. Для наглядности приведём частный случай одномерного переноса по~$x_1$ при~$U_1 > 0$ и отсутствии потоков по~$x_2$, когда баланс упрощается и даёт явную рекуррентную связь:
\begin{equation}\label{eq:jfo_theta_update}
  \vartheta_{i,j} = \frac{\vartheta_{i-1,j}\,h_{i-\frac{1}{2},j}}
    {h_{i+\frac{1}{2},j}}.
\end{equation}

Это соотношение выражает сохранение массы при одномерном переносе: входящий куэттовский поток через грань~$i-1/2$ равен выходящему потоку через грань~$i+1/2$. В двумерном случае баланс~\eqref{eq:jfo_couette_balance} включает потоки по обоим направлениям и реализуется в виде дискретной схемы переноса для~$\vartheta$ (раздел~3.2).

Для нестационарного случая в узлах, отнесённых к кавитационной зоне, используется явная по времени схема (временной шаг~$\Delta t$ из~\eqref{eq:discrete_balance}):
\begin{equation}\label{eq:jfo_theta_transient}
\begin{split}
  &(\vartheta\,h)_{i,j}^{n+1} = (\vartheta\,h)_{i,j}^{n} - {} \\
  &- \Delta t \left[
    \frac{F_{C,\,i+\frac{1}{2},j} - F_{C,\,i-\frac{1}{2},j}}{\Delta x_1}
    + \frac{G_{C,\,i,j+\frac{1}{2}} - G_{C,\,i,j-\frac{1}{2}}}{\Delta x_2}
  \right].
\end{split}
\end{equation}

После обновления выполняется отсечка:
\[
  \vartheta_{i,j} := \min\bigl(1,\;\max(0,\;\vartheta_{i,j})\bigr).
\]

Отсечка обеспечивает физичность решения ($0 \leq \vartheta \leq 1$). Если~$\vartheta_{i,j}$ после обновления достигает единицы, узел переходит в зону полного заполнения на следующей внешней итерации.

\textit{Обновление активного множества.}
После решения подсистемы для давления и проекции $p \leftarrow \max(p,\, p_{\mathrm{cav}})$ каждый узел классифицируется: если $p_{i,j} > p_{\mathrm{cav}} + \varepsilon$, узел принадлежит зоне полного заполнения и~$\vartheta_{i,j} = 1$; если $p_{i,j} \leq p_{\mathrm{cav}} + \varepsilon$, узел принадлежит зоне кавитации и~$\vartheta_{i,j}$ обновляется из баланса переносных потоков~\eqref{eq:jfo_couette_balance}. Малый параметр~$\varepsilon$ (порядка $10^{-6}$-$10^{-8}$ в безразмерном давлении или относительно характерного масштаба давления) необходим для устойчивого определения зоны кавитации и предотвращения осцилляций активного множества на границе зон.

Переключение узлов между зонами~-- это и есть обновление активного множества. Узел может переходить из кавитации в полное заполнение (если~$\vartheta$ достигает~1 или давление вырастает выше~$p_{\mathrm{cav}}$) и обратно (если после проекции выполняется~$p_{i,j} \leq p_{\mathrm{cav}} + \varepsilon$). Процесс продолжается до стабилизации.

\textit{Двухуровневая итерационная структура.}
Внутренний уровень~-- решение линейной системы для давления методом Red-Black SOR на GPU (раздел~3.4) при фиксированном~$\vartheta$. Сходимость внутреннего цикла контролируется по критерию~\eqref{eq:convergence_change} с параметрами из таблицы~\ref{tab:solver_params}. Внешний уровень~-- итерации по активному множеству. На каждой внешней итерации последовательно выполняются: решение для давления (внутренний цикл), проекция~$p \geq p_{\mathrm{cav}}$, обновление активного множества, обновление~$\vartheta$ в кавитационной зоне, отсечка. Общая схема алгоритма приведена в подразделе~2.3.3 (см.\ также рисунок~\ref{fig:jfo_algorithm}).

\textit{Критерии остановки внешнего цикла.}
Внешние итерации прекращаются при одновременном выполнении двух условий: стабилизации активного множества (ни один узел не переключился между зонами на текущей итерации) и малости изменения решения ($\max_{i,j} |p_{i,j}^{(m+1)} - p_{i,j}^{(m)}| < \varepsilon_p$, $\max_{i,j} |\vartheta_{i,j}^{(m+1)} - \vartheta_{i,j}^{(m)}| < \varepsilon_\vartheta$; критерии сформулированы в подразделе~2.3.3). На практике стабилизация активного множества достигается за 5-20~внешних итераций (в зависимости от задачи), при этом каждая внешняя итерация содержит полный внутренний цикл GPU~SOR.

\textit{Инициализация.}
Начальное приближение: $p^{(0)}_{i,j} = p_{\mathrm{cav}}$, $\vartheta^{(0)}_{i,j} = 1$ (полное заполнение). При параметрических расчётах используется тёплый старт: решение предыдущей задачи как начальное приближение для~$p$ и~$\vartheta$, что существенно сокращает число как внешних, так и внутренних итераций.

Описанная реализация при консервативной аппроксимации потоков обеспечивает массосохраняющее решение задачи кавитации и применяется ко всем рассматриваемым постановкам (упорный и радиальный подшипники, возвратно-поступательная пара). Результаты расчётов приведены в главах~5--7.

\textit{Вычислительная стоимость JFO-реализации.}
Основная вычислительная стоимость JFO-алгоритма складывается из~двух частей. Первая~--- внутренний GPU~SOR-цикл для~решения задачи на~давление~--- по~трудоёмкости аналогична обычному решению уравнения Рейнольдса и~масштабируется как~$O(N_1 N_2)$ на~итерацию. Вторая~--- обновление активного множества и~степени заполнения~$\vartheta$~--- выполняется за~один проход по~сетке и~требует значительно меньших затрат. Суммарная стоимость одного внешнего цикла JFO примерно в~2--3~раза превышает стоимость одного SOR-цикла простой Reynolds-задачи; общее число внешних итераций обычно составляет 10--30. Таким образом, JFO-расчёт оказывается примерно в~20--60~раз дороже, чем~один прогон Reynolds-решателя с~обнулением давления. На~практике это приемлемо для~параметрических расчётов при~использовании тёплого старта.

\textit{Особенности сходимости при~фрагментированной кавитации на~микротекстуре.}
При~текстурированных поверхностях с~большим числом лунок (порядка~50--200 на~расчётную область) кавитационная картина существенно фрагментируется: каждая лунка содержит небольшую собственную кавитационную зону. Это замедляет сходимость внешнего цикла активного множества по~следующей причине: при~инициализации с~полным заполнением алгоритму необходимо последовательно <<выявить>> все~локальные кавитационные участки, причём в~начале итераций граница одной кавитационной зоны может влиять на~соседние через~поле давления. Число внешних итераций в~этом случае может вырастать до~50--100 по~сравнению с~10--15 для~гладкого подшипника. Использование тёплого старта существенно снижает этот эффект: при~близких параметрах кавитационный рисунок от~предыдущего решения уже достаточно близок к~текущему, и~алгоритм сходится за~5--15 внешних итераций.

\textit{Допущение о~фиксированной кавитационной области в~динамике.}
При~расчёте динамических коэффициентов (раздел~4.4, глава~8) для~каждого набора возмущений $(\xi, \eta, \dot\xi, \dot\eta)$ в~принципе необходимо решать полную задачу JFO, в~том числе заново определять границу кавитации. Однако это многократно увеличивает вычислительные затраты. На~практике применяется допущение о~фиксированной кавитационной области (\textit{frozen cavitation}): граница кавитации берётся по~полю равновесного давления~$P_0$ и~при~решении возмущённой задачи не~пересматривается. Это допущение оправдано при~малых амплитудах возмущений~$|\xi|, |\eta| \ll c$ и~обеспечивает разумную точность для~прямых коэффициентов. Для~перекрёстных коэффициентов погрешность может быть заметнее, особенно в~режимах с~обширной кавитационной зоной.

\textit{Вычислительная стоимость JFO-реализации.}
Основная вычислительная стоимость JFO-алгоритма складывается из~двух частей. Первая~--- внутренний GPU~SOR-цикл для~решения задачи на~давление~--- по~трудоёмкости аналогична обычному решению уравнения Рейнольдса и~масштабируется как~$O(N_1 N_2)$ на~итерацию. Вторая~--- обновление активного множества и~степени заполнения~$\vartheta$~--- выполняется за~один проход по~сетке и~требует значительно меньших затрат. Суммарная стоимость одного внешнего цикла JFO примерно в~2--3~раза превышает стоимость одного SOR-цикла простой Reynolds-задачи; общее число внешних итераций обычно составляет 10--30. Таким образом, JFO-расчёт оказывается примерно в~20--60~раз дороже, чем~один прогон Reynolds-решателя с~обнулением давления. На~практике это приемлемо для~параметрических расчётов при~использовании тёплого старта.

\textit{Особенности сходимости при~фрагментированной кавитации на~микротекстуре.}
При~текстурированных поверхностях с~большим числом лунок (порядка~50--200 на~расчётную область) кавитационная картина существенно фрагментируется: каждая лунка содержит небольшую собственную кавитационную зону. Это замедляет сходимость внешнего цикла активного множества по~следующей причине: при~инициализации с~полным заполнением алгоритму необходимо последовательно <<выявить>> все~локальные кавитационные участки, причём в~начале итераций граница одной кавитационной зоны может влиять на~соседние через~поле давления. Число внешних итераций в~этом случае может вырастать до~50--100 по~сравнению с~10--15 для~гладкого подшипника. Использование тёплого старта существенно снижает этот эффект: при~близких параметрах кавитационный рисунок от~предыдущего решения уже достаточно близок к~текущему, и~алгоритм сходится за~5--15 внешних итераций.

\textit{Допущение о~фиксированной кавитационной области в~динамике.}
При~расчёте динамических коэффициентов (раздел~4.4, глава~8) для~каждого набора возмущений $(\xi, \eta, \dot\xi, \dot\eta)$ в~принципе необходимо решать полную задачу JFO, в~том числе заново определять границу кавитации. Однако это многократно увеличивает вычислительные затраты. На~практике применяется допущение о~фиксированной кавитационной области (\textit{frozen cavitation}): граница кавитации берётся по~полю равновесного давления~$P_0$ и~при~решении возмущённой задачи не~пересматривается. Это допущение оправдано при~малых амплитудах возмущений~$|\xi|, |\eta| \ll c$ и~обеспечивает разумную точность для~прямых коэффициентов. Для~перекрёстных коэффициентов погрешность может быть заметнее, особенно в~режимах с~обширной кавитационной зоной.
